Los capítulos anteriores han mostrado que cada técnica de generación tiene fortalezas y carencias complementarias. Las gramáticas formales producen salidas correctas y predecibles, pero rígidas y limitadas a las estructuras previstas. Los embeddings añaden información semántica pero no alteran la estructura de la generación. Los LLMs producen texto natural y flexible, pero sin garantías de fidelidad al contenido de entrada y con dependencia de infraestructura externa.

La hibridación parte de la hipótesis de que combinar estas técnicas puede producir resultados superiores a los de cada una por separado: una gramática que garantice la estructura, enriquecida con la naturalidad de un LLM, o un sistema que genere múltiples candidatos y seleccione el mejor. Se implementan tres estrategias de hibridación con objetivos diferenciados. Las dos primeras se basan en el patrón \textit{plan-then-realize}~\cite{moryossef2019stepbystep} y en la post-edición neuronal de texto generado por reglas~\cite{chen2025axsemantics}. La tercera implementa el paradigma de sobregeneración y ranking~\cite{sharma2019overgenerate}. Los resultados de la evaluación comparativa (capítulo~\ref{ch:evaluacion-comparativa}) permitirán determinar si esta hipótesis se confirma en la práctica.

\textit{(TODO: Añadir diagrama de arquitectura mostrando los tres pipelines y cómo reutilizan los motores existentes.)}

\section{Hibridación 1: Gramática + LLM como post-editor}
\label{sec:h-hybrid1}

Las Hibridaciones 1 y 2 se basan en dos patrones complementarios. El primero es el uso de un LLM como post-editor de sistemas de reglas, documentado por Chen et al.~\cite{chen2025axsemantics} en un sistema de producción de AX Semantics: el sistema de reglas garantiza la estructura y la fidelidad al contenido, mientras que el LLM aporta fluidez y naturalidad. El segundo es el patrón \textit{plan-then-realize}~\cite{moryossef2019stepbystep}, que separa la generación en un planificador simbólico que decide qué decir y un realizador neural que lo transforma en texto fluido. Elder et al.~\cite{elder2019intermediate} mostraron que esta separación reduce las alucinaciones respecto a enfoques puramente neurales.

La primera hibridación conecta la gramática formal con el LLM en un pipeline secuencial de dos pasos. Los pictogramas de entrada se procesan primero con la gramática, que intenta generar una frase estructurada. Si la gramática tiene éxito, la frase resultante se envía al LLM para su refinamiento. El LLM actúa como post-editor con ocho reglas estrictas que priorizan la fidelidad al contenido: no puede añadir ni eliminar conceptos, solo mejorar la concordancia, los artículos, el orden y la conjugación.

Si la gramática no reconoce la entrada, se activa un mecanismo de fallback: el LLM genera la frase desde cero. En una aplicación de CAA, dejar al usuario sin respuesta constituye un problema de accesibilidad más grave que devolver una frase imperfecta: la comunicación no puede quedar interrumpida porque el motor no reconozca un patrón.

Para la entrada ``quiero, agua'', la gramática detecta el patrón de necesidad simple y produce ``Quiero agua''. El post-editor evalúa que ya es correcta y la devuelve sin modificar. Para una entrada como ``ayer, abuelo, contar, cuento'', que no encaja en ningún patrón de la gramática, el LLM genera directamente ``Ayer el abuelo contó un cuento''.

Esta hibridación depende de un LLM externo para la post-edición, lo que introduce latencia de red, coste por llamada y falta de disponibilidad offline. La calidad de la post-edición depende enteramente de la ingeniería de prompts, y no existe un criterio formal para determinar si el conjunto de reglas del post-editor es óptimo ni completo.

\section{Hibridación 2: (Gramática + Embeddings) + LLM como post-editor}
\label{sec:h-hybrid2}

La segunda hibridación replica la estructura de H1 pero sustituye la gramática pura por el motor de embeddings, que combina gramática con embeddings semánticos, como fuente de la frase inicial. La frase generada por este motor, que ya incorpora refinamiento de preposiciones y ranking semántico, se envía al mismo post-editor LLM con las mismas ocho reglas.

La única diferencia con H1 es el motor que produce la frase intermedia: gramática pura en H1, gramática enriquecida con embeddings en H2. El post-editor es idéntico en ambos casos (mismo prompt, mismo modelo), lo que convierte la comparación en un experimento controlado: cualquier diferencia en los resultados se atribuye directamente a la contribución de los embeddings.

Las limitaciones son las mismas que las de H1: dependencia del LLM externo para la post-edición, con las implicaciones de latencia, coste y disponibilidad que ello conlleva.

\section{Hibridación 3: Overgenerate-and-Rank multimotor}
\label{sec:h-hybrid3}

La tercera hibridación se basa en el paradigma \textit{Generate-Filter-Rank}, formalizado por Sharma et al.~\cite{sharma2019overgenerate}: se generan múltiples candidatas, se filtran por criterios de calidad y se selecciona la mejor según una función de ranking. A diferencia de H1 y H2, que mejoran la salida de un motor individual, H3 busca la mejor frase entre todas las que los motores pueden producir. El pipeline se ejecuta en tres etapas secuenciales.

La primera etapa genera una candidata con cada uno de los cinco motores disponibles: gramática formal, embeddings, LLM, Hibridación 1 e Hibridación 2. Si algún motor falla, se descarta su candidata sin afectar al resto.

La segunda etapa evalúa la cobertura de pictogramas de cada candidata: se comprueba qué proporción de los pictogramas de entrada están representados en la frase generada, teniendo en cuenta las variaciones morfológicas del español. Solo pasan al ranking las candidatas que cubren el 100\% de los pictogramas. Si ninguna alcanza el umbral, todas pasan al ranking como fallback.

La tercera etapa calcula la perplejidad de cada candidata que pasó el filtro de fidelidad mediante un modelo de lenguaje en español~\cite{deepesp2019gpt2spanish} y las ordena de menor a mayor. La candidata con menor perplejidad, es decir, la que el modelo considera más natural, se selecciona como resultado final. Se aplica una normalización de puntuación terminal antes del cálculo, ya que la presencia o ausencia de un punto final altera significativamente la perplejidad sin reflejar diferencias reales de calidad.

H3 acumula la latencia de todos los motores que invoca, lo que la convierte en la opción más lenta. El paradigma de sobregeneración y ranking solo aporta valor cuando las candidatas son diversas y la función de ranking tiene granularidad suficiente para discriminar entre ellas; en un dominio con frases cortas y vocabulario restringido, las candidatas tienden a converger y el ranking pierde capacidad discriminativa. El umbral de fidelidad de cobertura completa puede descartar frases naturales que omiten un pictograma con representación implícita: ``Quiero agua'' para la entrada ``yo, quiero, agua'' podría fallar si el sistema no reconoce el sujeto implícito.
